\section{Problem-Level Ground-State Verification}
\label{sec:verification}
\subsection{Why reference bitstrings can misjudge a model}
QuantumQUBO Agent verifies small test cases by asking an LLM to translate human-stated solutions into reference
bitstrings in the generated formulation's variable layout\cite{mondal2026quantumqubo}. Enumeration accepts a
QUBO when its minimizers overlap with those references. The reference writer must count variables and populate
decision, slack and auxiliary bits without executing the model.

We audited all 33 failures in a local-model run with 67 successes on the 100-problem benchmark. Twelve failures
were shape mismatches. In at least seven, the generated test case had the wrong variable count; two arose from
model code and three remained undetermined. Six further failures concerned input-format mismatches and eleven
concerned argmin mismatches. A quadratic-assignment benchmark also stated an assignment costing 20 together
with an optimum of 12. Thus, failure against a generated reference need not indicate a wrong formulation.

\subsection{Verification in the language of the problem}
A case contains data $d$, the human-stated optimum $v^*$, an optional stated optimal solution and a shared
problem-level solution format. The generated builder includes a decoder. A separate Checker, written from the
problem statement without seeing the formulation, evaluates decoded decisions. Where a stated solution exists,
its evaluation must reproduce $v^*$ before the Checker is trusted on that case.

Let $G$ be the set of enumerated ground states. Acceptance requires
\begin{equation}
\begin{aligned}
 &\forall x\in G:\quad \mathrm{feasible}(\mathrm{decode}(d,x))\\
 &\qquad{}\land\ f(\mathrm{decode}(d,x))=v^*,
\end{aligned}
\end{equation}
with the implementation's numeric comparison tolerance. Checking every ground state rejects a degenerate QUBO
that also admits an invalid or non-optimal decision. Inconsistent reference/checker cases are marked
unverifiable rather than blamed on the model. Encodings beyond the enumeration threshold of 26 variables are
also unverifiable under this protocol. This exhaustive benchmark check is distinct from the sampled runtime
Requirement Checker in Section~\ref{sec:checkers}.

\subsection{Repeated benchmark results}
Table~\ref{tab:verification} retains unverifiable cases as a separate outcome. Five benchmarks failed under the
local bitstring protocol but succeeded in all three problem-level runs; three showed the reverse because their
ground states included non-optimal decisions. Of 16 unverifiable outcomes, ten exceeded the enumeration limit
and six reflected disagreement with benchmark references.

\begin{table}[t]
\centering\small
\caption{Results on 100 benchmarks. The original authors used remote models; local comparisons use the same locally served model.}
\label{tab:verification}
\begin{tabular}{@{}lrr@{}}\toprule
Protocol & Successes & Unverifiable \\
\midrule
Original authors' bitstring result & 68 & -- \\
Local bitstring run & 67 & -- \\
Problem-level, run 1 & 69 & 4 \\
Problem-level, run 2 & 72 & 6 \\
Problem-level, run 3 & 70 & 6 \\
\bottomrule\end{tabular}
\end{table}

The mean problem-level success rate is 70.3\% (range 69--72\%). Although temperature is zero, 27 benchmarks
changed outcome across repeated runs; 56 succeeded in all three and 17 in none. These observations support
repeated evaluation and recording per-problem outcomes. The protocols differ in both reference construction
and acceptance strictness, so their aggregate rates are not a clean causal estimate of one changed component.
